% e-TeX expression syntax: precedence and parentheses.
%
% Ours was a left-to-right fold, which is a different language:
% \numexpr 1+2*3 is 7, not 9. Every gate was green -- plain.tex never writes a
% mixed-precedence expression, and neither did any battery until this one.
% l3fp, which is built almost entirely out of parenthesised \__fp_int_eval:w,
% could not evaluate its first expression.
%
% The grammar is e-TeX's (etex.ch, scan_expr):
%   expr   := term { (+|-) term }
%   term   := factor { (*|/) integer-factor }
%   factor := operand | ( expr )
%
% The asymmetry in `term` is what makes a dimension expression type-check:
% 2pt*3 is a dimension times a *number*, so after the first factor the
% operands are integers however the expression as a whole is typed.
\input plain
\def\M#1{\message{<#1>}}
\count5=7 \dimen5=3pt \skip5=2pt plus 1pt
\M{a=\the\numexpr 1+2*3\relax}
\M{b=\the\numexpr (1+2)*3\relax}
\M{c=\the\numexpr 2*(3+4)-5\relax}
\M{d=\the\numexpr 100/(2*5)\relax}
\M{e=\the\numexpr ((1+2)*(3+4))\relax}
\M{f=\the\numexpr 1-2-3\relax}
\M{g=\the\numexpr 8/4/2\relax}
\M{h=\the\numexpr 2+3*4-5/5\relax}
\M{i=\the\numexpr 7/2\relax}
\M{j=\the\numexpr -7/2\relax}
\M{k=\the\numexpr 1+(2+(3+(4+5)))\relax}
\M{l=\the\numexpr \count5*2+1\relax}
\M{m=\the\numexpr (\count5+1)*2\relax}
\M{n=\the\numexpr 2*3/4\relax}
\M{o=\the\numexpr 32767*2+1\relax}
\M{p=\the\dimexpr (1pt+2pt)*3\relax}
\M{q=\the\dimexpr 1pt+2pt*3\relax}
\M{r=\the\dimexpr \dimen5*2+1pt\relax}
\M{s=\the\dimexpr (1pt+1pt)/2\relax}
\M{t=\the\dimexpr 1pt*(2+3)\relax}
\M{u=\the\glueexpr (1pt plus 1pt)*2\relax}
\M{v=\the\glueexpr \skip5*2+1pt plus 1pt\relax}
\M{w=\the\glueexpr (\skip5+1pt plus 1pt)*2\relax}
\M{x=\the\numexpr 1 + 2 * 3 \relax}
\M{y=\the\numexpr ( 1 + 2 ) * 3 \relax}
\bye
